\chapter{Azioni necessarie su GitHub}
\label{chap:github}

\section{Obiettivo}
In questa sezione vengono descritte le operazioni fondamentali da eseguire su GitHub per lavorare in modo ordinato sul progetto.

\section{Flusso operativo}
\begin{enumerate}
  \item Clonare la repository.
  \item Creare un branch dedicato.
  \item Apportare modifiche locali.
  \item Eseguire commit chiari.
  \item Eseguire push.
  \item Aprire una pull request.
  \item Aggiornare il branch locale.
\end{enumerate}

\section{Comandi principali}
\begin{itemize}
  \item \texttt{git clone <url>}
  \item \texttt{git checkout -b <branch>}
  \item \texttt{git add .}
  \item \texttt{git commit -m "messaggio"}
  \item \texttt{git push}
  \item \texttt{git pull}
\end{itemize}

\section{Buone pratiche}
\begin{itemize}
  \item commit piccoli e coerenti;
  \item nomi di branch chiari;
  \item pull request ben descritte.
\end{itemize}